% !TEX root = ../main.tex

\chapter{绪论}

\section{引言}
随着云办公、远程教学、在线医疗、云游戏和沉浸式交互应用的快速发展，实时通信（Real-Time Communication，RTC）已经从单一的视频通话工具演化为数字基础设施的重要组成部分。在这类应用中，用户对传输质量的要求具有显著的多目标特征：一方面需要尽可能高的媒体码率以保障画面清晰度，另一方面又要求端到端时延、抖动和卡顿率维持在较低水平，以确保交互过程的实时性和连续性\cite{Wang2024Pudica}。因此，拥塞控制算法能否在复杂网络环境中及时、稳定地调节发送速率，直接决定了 RTC 系统的服务质量与用户体验。
\par 当前 WebRTC 体系中广泛采用的 Google Congestion Control（GCC）算法\cite{Carlucci2016GCC}，本质上属于基于启发式规则的拥塞控制方法。该类算法依赖预设阈值、延迟梯度和丢包反馈进行码率调整，在网络环境相对稳定时具有较好的工程可用性。然而，随着无线网络成为 RTC 的主要接入场景，网络链路的时变性、随机性和异构性不断增强，启发式算法逐渐暴露出适应性不足的问题\cite{Zhang2020WebRTCGCC}。在 Wi-Fi、蜂窝网络等环境中，带宽可能在短时间内剧烈波动，链路还可能出现非拥塞型丢包、抖动突增和瞬时测量异常。固定规则驱动的拥塞控制策略往往难以及时区分“真实拥塞”与“无线扰动”，从而导致速率调整滞后、过度振荡或带宽利用不足，最终影响媒体流的流畅度和稳定性。
\par 近年来，强化学习为拥塞控制提供了新的研究视角。与传统方法相比，强化学习能够从历史状态与反馈中学习复杂的非线性决策策略，并可通过奖励函数同时考虑吞吐、时延、丢包和平滑性等多维目标，更接近真实用户体验的优化需求\cite{Bentaleb2023BoB}。以 OnRL\cite{Zhang2020OnRL} 为代表的在线强化学习方法已在移动视频通话中展示了学习码率控制策略的可行性。然而，在线强化学习通常依赖与真实环境的持续交互来完成探索和更新，这意味着训练过程可能直接干扰正在进行的 RTC 会话，带来明显的试错成本和部署风险\cite{Zhang2021Loki}。同时，网络环境具有明显的非平稳特征，真实系统中的探索样本难以系统覆盖长尾场景，导致策略泛化能力受限。
\par 离线强化学习（Offline Reinforcement Learning，Offline RL）为上述问题提供了可行的新路径。其核心思想是在不与环境持续交互的条件下，直接利用历史轨迹数据完成策略学习。对于 RTC 场景而言，大量由 GCC 等行为策略产生的网络状态、控制动作和反馈结果，本身就是高价值的离线训练资源。近期研究如 Mowgli\cite{Agarwal2025Mowgli}、Tan 等人\cite{Tan2024Accurate}以及 Du 等人\cite{Du2025LLLS}的工作已初步验证了利用历史数据进行被动学习或离线训练可以有效提升带宽预测的准确性与码率控制的鲁棒性。通过对这些轨迹进行清洗、建模和策略优化，离线强化学习有望在不影响在线用户体验的前提下，实现对传统启发式拥塞控制策略的性能提升。
\par 在工程实践中，GCC\cite{Carlucci2016GCC} 已成为 WebRTC 生态中最具代表性的拥塞控制方案。其通过接收端延迟变化趋势与发送端丢包信息的联合反馈，动态调整目标码率，在较多常规场景下具备较好的稳定性和可部署性。除 GCC 外，NADA\cite{Zhu2020NADA} 和 Copa\cite{Arun2018Copa} 等基于延迟的拥塞控制算法也被广泛研究，Pantheon\cite{Yan2018Pantheon} 平台则为各类拥塞控制方案的对比评估提供了基准环境。围绕 GCC 的分析、改进与替代研究构成了当前 RTC 拥塞控制的重要方向。为了提升拥塞控制策略的环境适应能力，学术界开始引入机器学习和强化学习方法\cite{Zhang2020OnRL,Zhang2021Loki,Bentaleb2023BoB}。相关工作普遍将带宽、时延、丢包率、历史码率等统计量作为状态输入，通过神经网络学习从网络状态到码率动作的映射关系。一些研究已经表明，强化学习方法在特定网络仿真环境下能够取得优于传统启发式算法的性能表现。
\par 然而，现有研究仍存在两方面突出问题。第一，许多强化学习方法依赖在线交互训练，训练过程需要不断试探新的动作策略，其探索行为可能显著降低真实通话质量，因此难以直接用于生产环境。第二，已有不少工作主要依赖网络仿真平台（如 ns-3\cite{Riley2010NS3}、SparkRTC\cite{SparkRTC2023}）完成训练与验证。虽然仿真工具便于控制变量和大规模实验，但在无线信道扰动、终端处理延迟、系统调度和浏览器实现细节等方面，仿真环境与真实 WebRTC 系统之间仍存在明显的 Sim-to-Real 差距\cite{Pinyoanuntapong2021SimToReal}。策略在仿真中取得的效果未必能够稳定迁移到真实设备与真实链路之上。社区曾推出 OpenNetLab\cite{Eo2022OpenNetLab} 和 AlphaRTC\cite{AlphaRTC2021} 等开源实验平台试图弥合这一鸿沟，但目前公网测试节点已停止维护，难以继续支持真实网络实验。
\par 在此背景下，离线强化学习逐渐成为新的研究热点。与在线强化学习不同，离线强化学习直接基于既有轨迹数据进行策略学习，避免了在线探索带来的风险，特别适用于高可靠、低容错的实时通信系统。隐式 Q 学习（Implicit Q-Learning，IQL）\cite{Kostrikov2021IQL}、保守 Q 学习（Conservative Q-Learning，CQL）\cite{Kumar2020CQL}等算法在离线决策任务上表现出较强潜力。相比需要频繁访问环境或依赖复杂行为约束的方案，IQL 在训练稳定性、超参数敏感性和工程实现复杂度方面更具优势，因此逐渐被用于需要从历史次优策略中挖掘改进空间的任务中\cite{Tan2024Accurate}。
\par 总体来看，国内外在 RTC 拥塞控制研究方面已经形成了“启发式方法长期主导、强化学习方法持续探索、离线强化学习逐渐兴起”的发展态势。但结合真实 WebRTC 端到端平台、真实设备采集数据、离线强化学习训练和在线 A/B 验证的完整闭环研究仍相对不足，尤其缺少一套能够兼顾数据采集、模型部署和实验评估的工程化验证平台。本文正是在这一背景下展开研究，力图将离线强化学习从理论算法进一步推进到真实 RTC 系统的部署验证层面。


\section{本文研究主要内容}
针对现有 RTC 拥塞控制方法在真实动态网络环境中存在的适应性不足、在线学习风险高以及仿真验证可信度有限等问题，本文以 WebRTC 真实通信场景为对象，构建基于离线强化学习的拥塞控制优化与部署验证方案。研究工作围绕平台搭建、数据构建、模型训练、推理部署与效果评估五个环节展开，形成一条完整的端到端技术路线。
\par 第一，搭建 WebRTC 采集与验证平台。本文基于浏览器端 WebRTC 能力（参考 Hairpin\cite{Meng2024Hairpin} 和 AlphaRTC\cite{AlphaRTC2021} 的开源实现），构建了用于数据集采集和 A/B 对比验证的实验平台，实现媒体流建立、网络状态周期采样、轨迹保存及场景切换控制。平台分别支持基础数据采集链路和在线 A/B Test 链路，能够在两台物理设备之间完成视频通话测试，并利用 NetEm\cite{Hemminger2005NetEm} 在物理链路上叠加可控的带宽限制与延迟抖动，记录 RTT、抖动、丢包率、接收吞吐和 GCC 估计带宽等关键指标，为离线训练提供真实样本来源。
\par 第二，构建离线强化学习训练数据集。本文将采集得到的 WebRTC 轨迹转换为强化学习所需的状态—动作—奖励—下一状态样本。以一段固定长度的历史窗口作为状态表示\cite{Bentaleb2023BoB,Tan2024Accurate}，选取发送速率、接收速率、RTT、丢包率和抖动等特征构成输入向量，以当前估计带宽作为动作，并设计综合考虑吞吐收益、时延惩罚、丢包惩罚和动作平滑性的 QoE 奖励函数。针对原始数据中可能存在的异常零点、短时缺口和长断点问题，在数据预处理阶段引入清洗与插值机制，以提升训练样本质量和状态表达稳定性。
\par 第三，采用 IQL 算法完成离线策略训练。基于构建好的离线数据集，本文选用隐式 Q 学习算法\cite{Kostrikov2021IQL}训练拥塞控制策略网络。该方法无需在策略更新中显式查询分布外动作的价值，从而在一定程度上缓解离线强化学习中的分布偏移问题。训练完成后导出 Actor 模型及对应的状态归一化配置，用于后续在线推理部署。
\par 第四，设计并实现在线推理部署方案。为了兼顾联调便利性与实时控制性能，本文实现了两种推理方式：一种是基于 Python 服务的远端推理接口，支持 HTTP 与 WebSocket 调用；另一种是将训练好的 Actor 模型导出为 ONNX 格式，在浏览器侧完成本地推理。前者便于快速验证模型输入输出逻辑，后者则避免了远程请求带来的额外时延，更适合在实时控制回路中直接驱动 RTCRtpSender.setParameters() 完成 maxBitrate 调节。
\par 第五，建立基于 QoE 的 A/B 测试评估流程。本文将 GCC\cite{Carlucci2016GCC} 作为对照策略，将离线强化学习策略作为实验策略，在相同的 WebRTC 平台和网络场景下进行对比实验，并从平均接收带宽、RTT 高分位值、丢包率、抖动、码率变化平滑性及综合 QoE 分数等多个维度评估算法效果，验证所提方案在真实应用链路中的实际收益。
\par 综合而言，本文的技术路线可以概括为：首先在真实 WebRTC 平台中采集 GCC 轨迹；其次将原始日志处理为适合离线强化学习训练的标准样本；随后利用 IQL 学习码率控制策略；最后通过在线推理与 A/B 测试完成部署验证。与仅在仿真器中研究算法不同，本文强调真实设备、真实轨迹和真实部署闭环，这也是本研究的主要特色所在。

\section{本文研究意义}
本文的研究意义体现在学术价值与工程应用两个层面。
\par 在学术层面，本文将离线强化学习方法引入即时通信媒体流拥塞控制任务，探索了利用历史 GCC 轨迹进行策略优化的可行性。与在线强化学习相比，离线学习范式无需承受在线探索对通话质量的干扰，能够安全地从次优行为策略的历史数据中提取改进信号。本文选用的 IQL 算法\cite{Kostrikov2021IQL}通过期望回归机制避免对分布外动作的价值查询，在理论层面对缓解离线学习中的分布偏移问题提供了有效保障。此外，本文面向真实 WebRTC 场景而非纯仿真环境开展研究，所获得的实验结论具有更强的工程参考价值，有助于缩小算法验证与实际部署之间的差距\cite{Pinyoanuntapong2021SimToReal}。
\par 在工程应用层面，本文的贡献主要体现在三个方面。其一，构建了覆盖“采集—训练—部署—评估”的完整工程闭环平台，使研究不仅停留在算法层面，还能够落到浏览器侧与服务侧的实际控制路径中，为后续研究者提供了可复用的实验基础设施。其二，在部署层同时支持远端 Python 推理和浏览器本地 ONNX 推理两种方案，并通过实验对比揭示了实时控制系统中推理延迟与被控系统耦合关系的重要性——当推理请求与被测媒体流共享物理信道时，远端推理方案在带宽受限场景下存在原理性缺陷，这一工程发现对后续智能拥塞控制的架构设计具有启示意义。其三，本文建立的 A/B 测试评估框架采用多维度 QoE 指标体系，能够系统、定量地比较不同拥塞控制策略在真实网络扰动下的表现，为面向实际部署的智能拥塞控制算法设计提供了规范化的评估方法论。

\section{本章小结}
本章首先介绍了实时通信拥塞控制领域的研究背景，分析了以 GCC 为代表的启发式算法在复杂网络环境下面临的适应性挑战，以及在线强化学习方法在部署安全性方面的固有局限，进而引出离线强化学习作为新的可行研究路径。随后，本章综述了国内外在 RTC 拥塞控制领域从启发式方法到强化学习方法再到离线强化学习的发展脉络，指出当前研究在真实端到端闭环验证方面的不足。在此基础上，本章概述了本文的五项主要研究内容——平台搭建、数据构建、模型训练、推理部署与 A/B 测试评估，并阐述了本文在学术与工程两个层面的研究意义。后续章节将依次展开各部分的详细论述。
